\documentclass[11pt, letterpaper]{article}

% --- Idioma y codificación ---
\usepackage[spanish]{babel}
\usepackage[utf8]{inputenc}
\usepackage[T1]{fontenc}
\usepackage{lmodern}

% --- Matemáticas ---
\usepackage{amsmath, amssymb, amsthm}

% --- Formato ---
\usepackage[margin=2.5cm]{geometry}
\usepackage{parskip}
\usepackage{booktabs}
\usepackage{hyperref}
\hypersetup{colorlinks=true, linkcolor=blue!70!black, urlcolor=blue!70!black}

% --- Entornos de teoremas ---
\theoremstyle{definition}
\newtheorem{definition}{Definición}[section]
\theoremstyle{plain}
\newtheorem{theorem}{Teorema}[section]
\newtheorem{lemma}[theorem]{Lema}
\newtheorem{corollary}[theorem]{Corolario}
\theoremstyle{remark}
\newtheorem{remark}{Observación}[section]

\title{%
  T02 --- Demostración formal:\\[4pt]
  Caso promedio de quicksort
}
\author{Curso Linux--C--Rust --- B15 Estructuras de datos en C}
\date{}

\begin{document}
\maketitle
\tableofcontents

% ===================================================================
\section{Modelo}
% ===================================================================

Consideramos quicksort con pivote aleatorio uniforme (cada elemento
tiene probabilidad $1/n$ de ser elegido como pivote). La entrada es una
permutación uniforme de $\{1, 2, \ldots, n\}$ (valores distintos).

Sea $C(n)$ la variable aleatoria que cuenta el número total de
comparaciones. Queremos calcular $E[C(n)]$.

% ===================================================================
\section{Observación clave: qué pares se comparan}
% ===================================================================

En lugar de analizar la recurrencia directamente, contamos
\textbf{qué pares de elementos se comparan} a lo largo de toda la
ejecución.

\begin{remark}
Dos elementos se comparan \textbf{a lo sumo una vez} durante
quicksort. Cuando se comparan, uno de ellos es el pivote de alguna
llamada recursiva. Después de la partición, el pivote queda en su
posición final y nunca participa en comparaciones futuras.
\end{remark}

\begin{definition}
Sea $z_i$ el $i$-ésimo menor elemento de la entrada (el elemento de
rango~$i$ en el array ordenado). Definimos la variable indicadora:
\[
  X_{ij} = \begin{cases}
    1 & \text{si } z_i \text{ y } z_j \text{ se comparan durante la
        ejecución} \\
    0 & \text{si no}
  \end{cases}
\]
para $1 \leq i < j \leq n$.
\end{definition}

El número total de comparaciones es:
\[
  C(n) = \sum_{i=1}^{n-1} \sum_{j=i+1}^{n} X_{ij}
\]

Por linealidad de la esperanza:
\[
  E[C(n)] = \sum_{i=1}^{n-1} \sum_{j=i+1}^{n} E[X_{ij}]
          = \sum_{i=1}^{n-1} \sum_{j=i+1}^{n}
            \Pr[z_i \text{ y } z_j \text{ se comparan}]
\]

% ===================================================================
\section{Demostración 1:
  \texorpdfstring{$\Pr[z_i \text{ y } z_j \text{ se comparan}]
  = \frac{2}{j-i+1}$}{Pr[zi y zj se comparan] = 2/(j-i+1)}}
% ===================================================================

\begin{lemma}\label{lem:prob}
Para $1 \leq i < j \leq n$:
\[
  \Pr[X_{ij} = 1] = \frac{2}{j - i + 1}
\]
\end{lemma}

\begin{proof}
Consideremos el conjunto
$Z_{ij} = \{z_i, z_{i+1}, \ldots, z_j\}$, que tiene $j - i + 1$
elementos.

\textbf{Observación clave:} $z_i$ y $z_j$ se comparan si y solo si
uno de ellos es el \textbf{primer} elemento de $Z_{ij}$ elegido como
pivote en alguna llamada recursiva.

¿Por qué? Mientras ningún elemento de $Z_{ij}$ haya sido pivote,
todos los elementos de $Z_{ij}$ permanecen en el mismo subarreglo
(porque ningún pivote externo los separa: un pivote $z_k$ con $k < i$
pone a todos los $Z_{ij}$ a su derecha, y un pivote $z_k$ con $k > j$
pone a todos los $Z_{ij}$ a su izquierda).

Cuando el primer elemento de $Z_{ij}$ se elige como pivote:
\begin{itemize}
  \item Si es $z_i$ o $z_j$: ese pivote se compara con todos los demás
    elementos del subarreglo, incluyendo el otro extremo. $z_i$ y
    $z_j$ \textbf{se comparan}. \checkmark
  \item Si es $z_k$ con $i < k < j$: la partición pone $z_i$ en el
    lado izquierdo (pues $z_i < z_k$) y $z_j$ en el lado derecho
    (pues $z_j > z_k$). $z_i$ y $z_j$ quedan en subproblemas
    disjuntos y \textbf{nunca se comparan}. $\times$
\end{itemize}

Como el pivote es aleatorio uniforme, cada elemento de $Z_{ij}$ tiene
la misma probabilidad $1/(j - i + 1)$ de ser el primero elegido. Los
casos favorables son $z_i$ o $z_j$ (2 de $j - i + 1$):
\[
  \Pr[X_{ij} = 1] = \frac{2}{j - i + 1} \qedhere
\]
\end{proof}

% ===================================================================
\section{Demostración 2:
  \texorpdfstring{$E[C(n)] = 2(n+1)H_n - 4n$}{E[C(n)] = 2(n+1)Hn - 4n}}
% ===================================================================

\begin{theorem}\label{thm:expected}
El número esperado de comparaciones de quicksort con pivote aleatorio
es:
\[
  E[C(n)] = 2(n+1)H_n - 4n
\]
donde $H_n = \sum_{k=1}^{n} 1/k$ es el $n$-ésimo número armónico.
\end{theorem}

\begin{proof}
Sustituimos $\Pr[X_{ij} = 1] = 2/(j-i+1)$:
\[
  E[C(n)] = \sum_{i=1}^{n-1} \sum_{j=i+1}^{n} \frac{2}{j - i + 1}
\]

Hacemos el cambio de variable $d = j - i$ (la ``distancia'' entre los
rangos):
\[
  = \sum_{i=1}^{n-1} \sum_{d=1}^{n-i} \frac{2}{d + 1}
\]

Intercambiamos el orden de la suma. Para un valor fijo de $d$, $i$ va
de~1 a $n - d$ (pues $j = i + d \leq n$ implica $i \leq n - d$). Hay
$n - d$ términos:
\[
  = \sum_{d=1}^{n-1} (n - d) \cdot \frac{2}{d + 1}
  = 2 \sum_{d=1}^{n-1} \frac{n - d}{d + 1}
\]

Separamos:
\[
  = 2 \sum_{d=1}^{n-1} \frac{n}{d+1}
  - 2 \sum_{d=1}^{n-1} \frac{d}{d+1}
\]

\textbf{Primer sumando.} Con $k = d + 1$, $k$ va de~2 a $n$:
\[
  2n \sum_{k=2}^{n} \frac{1}{k} = 2n(H_n - 1)
\]

\textbf{Segundo sumando:}
\[
  2 \sum_{d=1}^{n-1} \frac{d}{d+1}
  = 2 \sum_{d=1}^{n-1} \left(1 - \frac{1}{d+1}\right)
  = 2\left[(n{-}1) - \sum_{k=2}^{n} \frac{1}{k}\right]
  = 2(n{-}1) - 2(H_n - 1)
\]

\textbf{Combinando:}
\begin{align*}
  E[C(n)] &= 2n(H_n - 1) - [2(n{-}1) - 2(H_n - 1)] \\
          &= 2nH_n - 2n - 2n + 2 + 2H_n - 2 \\
          &= 2(n+1)H_n - 4n \qedhere
\end{align*}
\end{proof}

% ===================================================================
\section{Demostración 3:
  \texorpdfstring{$E[C(n)] \approx 1.39 \cdot n \log_2 n$}%
  {E[C(n)] approx 1.39 n log2 n}}
% ===================================================================

\begin{theorem}\label{thm:asymptotic}
$E[C(n)] \sim 2n \ln n \approx 1.386 \cdot n \log_2 n$.
\end{theorem}

\begin{proof}
Usamos la asíntota del número armónico:
\[
  H_n = \ln n + \gamma + O(1/n)
\]
donde $\gamma \approx 0.5772$ es la constante de Euler--Mascheroni.

Sustituyendo en $E[C(n)] = 2(n+1)H_n - 4n$:
\begin{align*}
  E[C(n)] &= 2(n+1)(\ln n + \gamma + O(1/n)) - 4n \\
          &= 2n\ln n + 2\ln n + 2(n+1)\gamma + O(1) - 4n \\
          &= 2n\ln n + (2\gamma - 4)n + O(\ln n)
\end{align*}

El término dominante es $2n\ln n$. Convertimos a $\log_2$:
\[
  2n\ln n = 2n \cdot \frac{\log_2 n}{\log_2 e}
          = 2n \cdot \log_2 n \cdot \ln 2
          = 2\ln 2 \cdot n\log_2 n
\]

Como $2\ln 2 \approx 1.3863$:
\[
  E[C(n)] \approx 1.386 \cdot n \log_2 n
\]
\end{proof}

\begin{remark}[Comparación con merge sort]
\[
  \frac{E[C_{\text{quicksort}}(n)]}{C_{\text{merge,worst}}(n)}
  \approx \frac{1.386 \cdot n\log_2 n}{n\log_2 n} \approx 1.39
\]
Quicksort hace ${\sim}39\%$ más comparaciones que merge sort en el
peor caso de merge sort. Sin embargo, quicksort es más rápido en la
práctica porque sus comparaciones tienen mejor localidad de cache
(opera in-place sobre un segmento contiguo).
\end{remark}

% ===================================================================
\section{Demostración 4: peor caso
  \texorpdfstring{$\Theta(n^2)$}{Theta(n\^2)}}
% ===================================================================

\begin{theorem}\label{thm:worst}
Quicksort con pivote fijo (primer o último elemento) tiene peor caso
$\Theta(n^2)$.
\end{theorem}

\begin{proof}
Si la entrada está ordenada y el pivote es el último elemento, cada
partición produce un subarreglo de tamaño $n - 1$ y uno vacío:
\[
  T(n) = T(n-1) + cn
\]

Resolviendo por expansión:
\[
  T(n) = cn + c(n{-}1) + c(n{-}2) + \ldots + c
       = c\sum_{k=1}^{n} k
       = \frac{cn(n+1)}{2}
       = \Theta(n^2) \qedhere
\]
\end{proof}

% ===================================================================
\section{Resumen de resultados}
% ===================================================================

\begin{table}[h]
\centering
\begin{tabular}{llp{5.5cm}}
\toprule
\textbf{Resultado} & \textbf{Técnica} & \textbf{Clave} \\
\midrule
$\Pr[z_i, z_j \text{ comp.}] = \frac{2}{j{-}i{+}1}$
  & Primer pivote en $Z_{ij}$
  & Solo $z_i$ o $z_j$ causan comparación \\
$E[C(n)] = 2(n{+}1)H_n - 4n$
  & Linealidad de la esperanza
  & Doble suma $\to$ suma sobre distancias \\
$E[C(n)] \approx 1.386 \cdot n\log_2 n$
  & Asíntota $H_n \sim \ln n$
  & $2\ln 2 \approx 1.386$ \\
Peor caso $\Theta(n^2)$
  & Expansión de $T(n{-}1) + cn$
  & Partición maximalmente desbalanceada \\
\bottomrule
\end{tabular}
\end{table}

\end{document}
